% texte.tex — f163v, vv. 80–122 (conversion initiale des travaux antérieurs)
\vers[n=80]{}{}{ἐπείπερ ἐστὶ Θεσμοφορίων ἡ μέση.}{puisque nous sommes au jour intermédiaire des Thesmophories.}
\vers[n=81]{ΕΥ.}{Euripide}{τοῦτ᾽ αὐτὸ γὰρ τοι κἀπολεῖν με προσδοκῶ.}{C'est précisément ce qui fait que je m'attends à périr.}
\vers[n=82]{}{}{αἱ γὰρ γυναῖκες ἐπιβεβουλεύκασί μοι}{Car les femmes ont monté un coup contre moi.}
\vers[n=83]{}{}{κἀν Θεσμοφόροιν μέλλουσι περί μου τήμερον}{Elles doivent, dans le temple des Thesmophories, à mon sujet aujourd'hui,}
\vers[n=84]{}{}{ἐκκλησιάζειν ἐπ᾽ ὀλέθρῳ.}{tenir assemblée afin de me perdre.}
\vers[suite]{ΜΝ.}{Mnésiloque}{τιὴ τί δή;}{Pourquoi cela ?}
\vers[n=85]{ΕΥ.}{Euripide}{ὁτιὴ τραγῳδῶ καὶ κακῶς αὐτὰς λέγω.}{Parce que je les théâtralise et parle mal d'elles.}
\vers[n=86]{ΜΝ.}{Mnésiloque}{νὴ τὸν Ποσειδῶ καὶ δικαίαν ἂν πάθοις.}{Par Poseidôn, tu serais châtié en justice.}
\vers[n=87]{}{}{ἀτὰρ τίν᾽ ἐκ ταύτης σὺ μηχανὴν ἔχεις;}{Mais, pour y échapper, de quelle machine disposes-tu ?}
\vers[n=88]{ΕΥ.}{Euripide}{Ἀγάθωνα πεῖσαι τὸν τραγῳδοδιδάσκαλον}{De persuader à Agathon, le poète tragique,}
\vers[n=89]{}{}{ἐς Θεσμοφόροιν ἐλθεῖν.}{de se rendre aux Thesmophories.}
\vers[suite]{ΜΝ.}{Mnésiloque}{τί δράσοντ᾽; εἰπέ μοι.}{Pour quoi faire, dis-moi ?}
\vers[n=90]{ΕΥ.}{Euripide}{ἐκκλησιάσοντ᾽ ἐν ταῖς γυναιξὶ χἂν δέῃ}{Pour intervenir à l'assemblée devant les femmes et, autant qu'il est nécessaire,}
\vers[n=91]{}{}{λέξονθ᾽ ὑπὲρ ἐμοῦ.}{me défendre.}
\vers[suite]{ΜΝ.}{Mnésiloque}{πότερα φανερὸν ἢ λάθρᾳ;}{Ouvertement ou en tapinois ?}
\vers[n=92]{ΕΥ.}{Euripide}{λάθρᾳ, στολὴν γυναικὸς ἠμφιεσμένον.}{En tapinois, revêtu d'un costume féminin.}
\vers[n=93]{ΜΝ.}{Mnésiloque}{τὸ πρᾶγμα κομψὸν καὶ σφόδρ᾽ ἐκ τοῦ σοῦ τρόπου·}{C'est ingénieux et bien dans tes façons.}
\vers[n=94]{}{}{τοῦ γὰρ τεχνάζειν ἡμέτερος ὁ πυραμοῦς.}{En fait d'artifices, à nous le gâteau.}
\vers[n=95]{ΕΥ.}{Euripide}{σίγα.}{Silence.}
\vers[suite]{ΜΝ.}{Mnésiloque}{τί δ᾽ ἔστιν;}{Qu'y a-t-il ?}
\vers[suite]{ΕΥ.}{Euripide}{Ἁγάθων ἐξέρχεται.}{Agathon sort.}
\vers[n=96]{ΜΝ.}{Mnésiloque}{καὶ ποῖός ἐστιν;}{Quel homme est-ce ?}
\vers[suite]{ΕΥ.}{Euripide}{οὗτος οὑκκυκλούμενος.}{Celui-là qu'on roule dehors.}
\vers[n=97]{ΜΝ.}{Mnésiloque}{ἀλλ᾽ ἦ τυφλὸς μέν εἰμ᾽· ἐγὼ γὰρ οὐχ ὁρῶ}{Est-ce donc que je suis aveugle ? Je ne vois pas qu'il y ait}
\vers[n=98]{}{}{ἄνδρ᾽ οὐδέν᾽ ἐνθάδ᾽ ὄντα, Κυρήνην δ᾽ ὁρῶ.}{un homme ici. C'est Cyrène que je vois.}
\vers[n=99]{ΕΥ.}{Euripide}{σίγα· μελῳδεῖν νῦν παρασκευάζεται.}{Chut ! Il se dispose à chanter.}
\vers[n=100]{ΜΝ.}{Mnésiloque}{Μύρμηκος ἀτραπούς, ἢ τί διαμινύρεται;}{Des sentes de Fourmi, ou que sont ces fredons ?}
\vers[n=101, indent=4em, findent=2.4em, metre={lyr.}]{ΑΓ.}{Agathon}{ἱερὰν Χθονίαιν δε-}{1° — Des deux Chtoniennes, recevez}
\vers[n=101, indent=4em, findent=2.4em]{}{}{ξάμεναι λαμπάδα κοῦραι}{la torche consacrée, vierges,}
\vers[n=102, indent=4em, findent=2.4em]{}{}{ξὺν ἐλευθέρᾳ πατρί-}{et avec votre libre patrie}
\vers[suite]{ΑΓ.}{Agathon}{\alinea δι χορεύσασθε βοάν.}{entonnez un chœur de danse.}
\vers[n=104, indent=4em, findent=2.4em]{}{}{τίνι δαιμόνων ὁ κῶμος;}{2° — À quelle divinité ce cômos ?}
\vers[n=105, indent=4em, findent=2.4em]{}{}{λέγε νιν εὐπείστως δὲ τοὐμὸν}{Parle. Je me laisse aisément induire}
\vers[n=106, indent=4em, findent=2.4em]{}{}{δαίμονας ἔχει σεβίσαι.}{à honorer les dieux.}
\vers[n=107, indent=4em, findent=2.4em]{}{}{ἄγε νῦν ὄλβιζε Μοῦσα}{3° — Allons, Muse, glorifie}
\vers[n=108, indent=4em, findent=2.4em]{}{}{χρυσέων ῥύτορα τόξων}{l'archer à l'arc d'or,}
\vers[n=109, indent=4em, findent=2.4em]{}{}{Φοῖβον, ὃς ἱδρύσατο χώρας}{Phébus, qui fonda les murs}
\vers[n=110, indent=4em, findent=2.4em]{}{}{γύαλα Σιμουντίδι γᾷ.}{de cette cité dans la contrée du Simois.}
\vers[n=111, indent=4em, findent=2.4em]{}{}{χαῖρε καλλίσταις ἀοιδαῖς}{4° — Sois salué de nos plus beaux chants,}
\vers[n=112, indent=4em, findent=2.4em]{}{}{Φοῖβ᾽ ἐν εὐμούσοισι τιμαῖς}{Phébus, qui dans les fêtes des Muses}
\vers[n=113, indent=4em, findent=2.4em]{}{}{γέρας ἱερὸν προφέρων.}{emportes le prix sacré.}
\vers[n=114, indent=4em, findent=2.4em]{}{}{τάν τ᾽ ἐν ὄρεσι δρυογόνοισι}{5° — Et chantez la Vierge qui se plaît dans les}
\vers[n=115, indent=4em, findent=2.4em]{}{}{κόραν ἀείσατ᾽ Ἄρτεμιν ἀγροτέραν.}{montagnes dryogènes, Artémis-chasseresse.}
\vers[n=117, indent=4em, findent=2.4em]{}{}{ἕπομαι κλῄζουσα σεμνὰν}{6° — À ta suite je célèbre et glorifie}
\vers[n=118, indent=4em, findent=2.4em]{}{}{γόνον ὀλβίζουσα Λατοῦς}{l'auguste rejeton de Léto,}
\vers[n=119, indent=4em, findent=2.4em]{}{}{Ἄρτεμιν ἀπειρολεχῆ.}{Artémis ignorant l'hymen.}
\vers[n=120, indent=4em, findent=2.4em]{}{}{Λατώ τε κρούματά τ' Ἀσιάδος ποδὶ}{7° — Ainsi que Léto et les percussions de}
\vers[n=121, indent=4em, findent=2.4em]{}{}{παράρυθμ' εὔρυθμα Φρυγίων}{la cithare d'Asie, aux beaux rythmes contrariés marqués par le pied}
\vers[n=122, indent=4em, findent=2.4em]{}{}{διὰ νεύματα Χαρίτων.}{sous la direction des Charites phrygiennes.}
